\section{Scriptural Date and Location}
{\small Each passage this Mass reads or sings from Scripture, once, in
canonical order; beneath each, its traditional frame and the modern critical
limit warranted by the stated sources. The Collect, Secret and Postcommunion
are composed liturgical prayers and carry no biblical date. Alternatives and
unresolved states remain explicit.\par}

\begin{dossiertable}
Gradual & Ps.~32:6, 12 (modern~33) & No place named in the appointed verses &
\chronodate{gradual}{\chronologyannotation{gradual}}\\
\cmidrule(lr){1-4}
\dossierprose{The appointed verses join two of the psalm's movements: the
blessedness of the people the Lord chose for his inheritance (v.~12) and the
heavens established by his word and spirit (v.~6). The Douay--Rheims heading
reads ``A psalm for David'', an attribution without a named occasion or place
of writing, first audience or stage in the attributed author's life. No
traditional era is established for this psalm. The NABRE introduction supplies
the broad limit shown for the Psalter, not a precise date for this poem, and
denies that individual psalms can be dated with certainty.}
\midrule
Communion & Ps.~75:12--13 (modern~76) & Judea and Sion, named earlier in the
psalm as the place of God's dwelling and victory (vv.~2--3) &
\chronodate{communion}{\chronologyannotation{communion}}\\
\cmidrule(lr){1-4}
\dossierprose{The superscription ascribes the psalm to Asaph, ``a canticle to
the Assyrians'' (v.~1), without attaching a date or era. In the complete
psalm God breaks the weapons of war in Sion and rises to judgment ``to save
all the meek of the earth'' (vv.~4, 10); the appointed conclusion commands
vows and their payment and gathers gift-bearers before the Lord who humbles
the spirit of princes. Neither title nor verses fixes an occasion, first
audience, stage in the attributed author's life or place of composition. The
Date column therefore carries only the Psalter-wide modern limit; this
individual psalm is not securely dated.}
\midrule
Alleluia & Ps.~101:2 (modern~102) & Sion named in the psalm's turn from
lament to hope (vv.~14, 17) &
\chronodate{alleluia}{\chronologyannotation{alleluia}}\\
\cmidrule(lr){1-4}
\dossierprose{The superscription describes the poor man's anxious prayer
and supplication before the Lord (v.~1): an afflicted speaker, no author,
no traditional date, first audience or authorial life-stage. The appointed
verse is the psalm's opening cry for a hearing; the complete poem moves from
that frailty to the rebuilding of Sion and the gathering of peoples to serve
the Lord. Sion is the object of the psalm's hope, not an established place of
writing, and the general Psalter boundary does not date any particular
restoration.}
\midrule
Introit & Ps.~118:1, 124, 137 (modern~119) & No place named &
\chronodate{introit}{\chronologyannotation{introit}}\\
\cmidrule(lr){1-4}
\dossierprose{An alphabetical psalm in praise of the law, with no named
author and no traditional date or setting. The antiphon places the confession
of God's just judgment before the petition for mercy (vv.~137, 124), reversing
their order in the complete psalm, and the verse sung with it is the psalm's
opening beatitude over those who walk in the law of the Lord (v.~1). No
particular crisis, first audience or writing-place is established by the
appointed verses, nor can an authorial life-stage be established where no
author is named. The date shown is only the Psalter-wide modern limit; this
individual psalm is not securely dated.}
\midrule
Offertory & Dan.~9:17--19, adapted & Daniel at the court of Babylon, in the
first year of Darius the Mede (9:1); prayer for the desolate sanctuary of
Jerusalem (9:17) &
\chronodate{offertory}{\chronologyannotation{offertory}}\\
\cmidrule(lr){1-4}
\dossierprose{Traditional attribution: the prophet Daniel, in the Exile,
recording his own prayer (``I, Daniel'', 9:2) when he understood from
Jeremias that the seventy years of desolation were ending. The 1908
\work{Catholic Encyclopedia} article that admits 570--536~B.C.\ as the book's
date grounds the traditional position in details it says only ``a resident in
Babylon'' could supply, written by Daniel, ``owing to his position at the
court of Babylon'', ``for the comfort of the Jews of his time and of
subsequent ages''. Gigot also reports the critical attribution to an unknown
later author rather than Daniel; his account describes scholarship of 1908,
not a fresh survey of present opinion, and it establishes no numerical modern
critical date for that alternative. The ``seventy years'' carried in the Date column is the
duration of the servitude within the prayer's setting (Jer.~25:11--12), not a
composition date. The chant selects and rearranges Daniel's petitions; its
liturgical form has its own history, distinct from the composition of the
chapter.}
\midrule
Gospel & Matt.~22:34--46 & Written, in the received account, in Judea before
Matthew's departure from Jerusalem, for the Jews who had believed; the
encounter narrated in the Temple at Jerusalem (21:23) &
\chronodate{gospel}{\chronologyannotation{gospel}}\\
\cmidrule(lr){1-4}
\dossierevent{Narrated event: the last of the Lord's controversies in the
Temple (21:23--22:46), shortly before the Passion forecast at 26:1--2. No
numerical narrated-event date for these verses is established here.}
\cmidrule(lr){1-4}
\dossierprose{Traditional attribution: St~Matthew the Apostle, whom the
Missal's heading names (\latin{secundum Matthaeum}). The received frame comes
from the early writers as the 1911 \work{Catholic Encyclopedia} article
reports them: the ecclesiastical writers agree that Matthew wrote for the
Jews, and Jerome summarizes the tradition that he published his Gospel in
Judea, in the Hebrew tongue, chiefly for those among the Jews who had
believed in Jesus. The Date column keeps the article's disputed alternatives,
which are not of equal standing. About A.D.~38--45 rests on early writers
counting eight years (Eusebius in his Chronicle, Theophylact, Euthymius
Zigabenus) or fifteen years (Nicephorus Callistus) from the Ascension. About
A.D.~40--42 follows the tradition that the apostles separated twelve years after
the Ascension, a tradition the article itself calls ``admittedly not too
reliable''. The years A.D.~40--45 report the article's survey of Catholic critics
of its own day, who in general favoured that range. About A.D.~60--68 is the
conditional alternative that results if, following Eusebius, the definitive
departure of the apostles is fixed later. About A.D.~64--67 rests on a
text of Irenaeus---Matthew writing while Peter and Paul were evangelizing and
founding the Church of Rome---whose difficulties of interpretation, the
article says, prevent any positive conclusion. About A.D.~50 is Durand's
date in the 1912 article on the New Testament, given for the Aramaic original
alone; he holds that nothing definite is known of when it was rendered into
Greek. These alternatives concern composition, while the Temple locates the
narrated controversy; no numerical date for that event is established here.
A separately generated comparison records the NABRE's post-A.D.~70 boundary
for the Greek Gospel, probably at least a decade later. It supplies no precise
endpoint and remains distinct from the default traditional alternatives.}
\midrule
Epistle & Eph.~4:1--6 & Paul, ``a prisoner in the Lord'' (4:1), writing in
captivity; to the saints at Ephesus (1:1) &
\chronodate{epistle}{\chronologyannotation{epistle}}\\
\cmidrule(lr){1-4}
\dossierprose{Traditional attribution: St~Paul, in one of the letters of his
captivity. The 1909 \work{Catholic Encyclopedia} article assigns the
captivity letters to the period A.D.~58--63 shown in the Date column,
judging that they follow the Epistle to the Romans and precede the first
persecution, and it leaves the place---Caesarea or Rome---``a much mooted
question'', while remarking that the liberty and
activity the letters display suit the Roman captivity better. The
encyclopedia's 1911 chronology of St~Paul assigns Ephesians, with Philemon,
Colossians and Philippians, the single year shown. The first hearers are the
church addressed in the received text of 1:1, gentiles once ``afar off'' and
now ``made nigh'' (2:11--18). The same 1909 article reports that most liberal
critics since Schleiermacher have denied the letter to Paul, a position for
which it gives no numerical date; no additional modern critical date is
established here. It also holds the words ``at Ephesus'' in 1:1 to be no part
of the primitive text and reads the letter as a circular to churches Paul had
not visited, probably in Asia Minor. The address above is the
received biblical text, not a settled identification of the original
recipients.}
\end{dossiertable}
